% from ega0-21-fr.tex
\oldpage[0\textsubscript{IV-1}]{161}
De sorte que l'on a la suite exacte
\begin{equation}
\tag{21.3.2.5}\label{0.21.3.2.5-fr}
0\to\Xi_{B/A}\to\Theta_{B/A}\xrightarrow{\pi_{B/A}}\Omega^1_A
\end{equation}
(on remarquera que ces notations peuvent prêter à confusion puisque
$\Theta_{B/A}$ et $\Xi_{B/A}$ dépendent non seulement de $A$ et de $B$,
mais aussi de $i$ et de $j$).
\end{env}

\begin{env}[21.3.3]
\label{0.21.3.3-fr}
Comme, en vertu de (21.3.1.2), on a $F_B=i\circ j$, on peut écrire pour
tout $B$-module $M$ (cf. (21.1.4))
\begin{equation}
\tag{21.3.3.1}\label{0.21.3.3.1-fr}
M^{(p)}=M\otimes_B B^{(p)}=(M\otimes_B A_{[j]})\otimes_A B_{[i]},
\end{equation}
d'où en particulier
\begin{equation}
\tag{21.3.3.2}\label{0.21.3.3.2-fr}
(\Omega^1_{B/A})^{(p)}=\Theta_{B/A}\otimes_A B_{[i]},
\end{equation}
et l'on déduit donc de (21.3.2.2) un homomorphisme canonique de $B$-modules
\begin{equation}
\tag{21.3.3.3}\label{0.21.3.3.3-fr}
\pi_{B/A}\otimes1_B:(\Omega^1_{B/A})^{(p)}\longrightarrow
\Omega^1_A\otimes_A B_{[i]}.
\end{equation}
Compte tenu de (20.5.4.2), l'image de cet homomorphisme est le $B$-module
engendré par les éléments $d_A(j(b))\otimes1$ pour $b\in B$. Leur image
par l'homomorphisme canonique $i_{B/A}:\Omega^1_A\otimes_A B\to\Omega^1_B$
déduit de $i$ (20.5.2) est donc dans le sous-$B$-module de $\Omega^1_B$
engendré par les $d_B(i(j(b)))$ pour $b\in B$~; en vertu de (21.3.1.2),
cette image est nulle. Autrement dit, on a une suite d'homomorphismes
\begin{equation}
\tag{21.3.3.4}\label{0.21.3.3.4-fr}
0\to\Xi_{B/A}\otimes_A B_{[i]}\to(\Omega^1_{B/A})^{(p)}
\to\Upsilon_{B/A}
\end{equation}
qui n'est pas nécessairement exacte, mais où le composé de deux
homomorphismes consécutifs est nul.
\end{env}

\begin{proposition}[21.3.4]
\label{0.21.3.4-fr}
Si $B$ est un $A$-module plat (pour $i$), la suite (21.3.3.4) est exacte.
\end{proposition}
\begin{proof}
Cela résulte de la définition de la platitude. La proposition s'applique en
particulier lorsque $B^p\subset A\subset B$, $i$ et $j$ étant respectivement
l'injection canonique et $F_B$, et que $B$ admet une $p$-base sur $A$ (de sorte
que $B$ est alors un $A$-module libre). Mais même dans ce cas le noyau
$\Xi_{B/A}$ n'est pas nécessairement nul. Toutefois~:
\end{proof}

\begin{proposition}[21.3.5]
\label{0.21.3.5-fr}
Soient $B$ un anneau réduit, $A$ un sous-anneau de $B$ tel que
$B^p\subset A\subset B$. Supposons qu'il existe une $p$-base
$(x_\lambda)_{\lambda\in L}$ de $B$ sur $A$ et une $p$-base
$(y_\mu)_{\mu\in M}$ de $A$ sur $B^p$~; alors l'homomorphisme canonique
(21.3.3.4)
\begin{equation}
\tag{21.3.5.1}\label{0.21.3.5.1-fr}
(\Omega^1_{B/A})^{(p)}\longrightarrow\Upsilon_{B/A}
\end{equation}
est bijectif, et les éléments $d_A(x_\lambda^p)\otimes1$ forment une base
du $B$-module $\Upsilon_{B/A}$.
\end{proposition}

\begin{proof}
